\documentclass[11pt]{article}

\usepackage[english]{babel}
\usepackage[latin1]{inputenc}
\usepackage{graphicx}
\usepackage{amssymb}
\usepackage{url}
\usepackage{enumerate}
\usepackage[nodayofweek]{datetime}
\usepackage{xcolor}
\usepackage{natbib}
\newcommand{\Real}{\mathbb{R}}


\textwidth 15cm
\setlength{\textheight}{1.1\textheight}
\newcommand\hi{\hspace*{\parindent}}
\newcommand\vi{\vspace{\baselineskip}}
\newcommand\lac{{\mbox{{\Huge\bf L}\hspace{-0.65em}
\raisebox{-1.2ex}{\Huge\bf A}\hspace{-1.1em}
\raisebox{-0.6ex}{\Huge\bf C~}}}}

\usepackage{fancyhdr}
\pagestyle{fancy}
\fancyhf{}
\lhead{Manuscript Reference Number: SCICO-D-17-00062}
\rhead{Page \thepage\ of \pageref{LastPage}}
\lfoot{}
\rfoot{}

\pagenumbering{roman}
\begin{document}
\setcounter{page}{1}
\thispagestyle{empty}

\hoffset -1.5cm \voffset .5cm

\sf

\vspace*{-2.5cm} \hspace*{-0.1cm} {\small
{\mbox{\begin{minipage}{1.5cm}
\centerline{\includegraphics[width=1.8cm]{logo_ufam.png}}
\end{minipage}}} \hspace*{0.5cm}
\begin{minipage}{11.0cm}{\sc \large \textbf{FEDERAL UNIVERSITY OF AMAZONAS}}\\
{\sc Faculty of Technology\\}
{\sc Department of Electronics and Computing}
\end{minipage}
} \vspace*{0mm}

\hspace*{-.7cm} {\rule[-1ex]{15cm}{0.03cm}}

\vi \vi
\begin{flushright}
\begin{minipage}{7.0cm}\small
{\bf Please Reply to:}\\
Mr. Felipe R. Monteiro\\
Federal University of Amazonas\\
Department of Electronics and Computing\\
Avenida General Rodrigo Ot\'avio, 6200, Coroado\\
69077-000 - Manaus - AM - BRAZIL\\
{\em e-mail: felipemonteiro@ufam.edu.br}
\end{minipage}
\end{flushright}

\vi \vi
\hspace*{\fill}{\small Manaus, \today.}
\vi \vi

\hoffset -1.5cm \voffset .5cm

\begin{flushleft}
Dr. Mark van den Brand\\
   Editor-in-Chief \\
   Science of Computer Programming Software Track
    \end{flushleft}
\vi

\begin{flushleft}
\textbf{REF.:} SCICO-D-17-00062
    \end{flushleft}
\vi \vi
 
\indent Dear Dr. Mark van den Brand,
\vi 

We thank you for asking us to comply with the Referees' reports, which allowed us to submit a revised version of our paper, entitled as ``{\em ESBMC-GPU A Context-Bounded Model Checking Tool to Verify CUDA Programs},'' by Felipe R. Monteiro, Erickson H. da S. Alves, Isabela S. Silva, Hussama A. Ismail, Lucas C. Cordeiro, and Eddie B. de Lima Filho.


The details of all changes made to our paper and also our response to each comment raised by the reviewers have been included at the end of this letter.

\vi\vi

\indent
~~~~~~~~~~~~~~~~~~~~~~~~~~~Sincerely,\\



\begin{quote}
\begin{quote}
\begin{flushright}


\vi
\vi

Felipe R. Monteiro~~~~~~~
\end{flushright}
\end{quote}
\end{quote}

\hoffset -1.5cm \voffset .5cm


% =============================================
\newpage
\subsection*{Authors' Response to the Review Comments on manuscript ``{ESBMC-GPU A Context-Bounded Model Checking Tool to Verify CUDA Programs}'' -- Manuscript Reference Number: SCICO-D-17-00062}

\vi

The authors thank the anonymous reviewers for their valuable and constructive comments, which helped us significantly improve the quality of our manuscript. 

\vi

Our manuscript has been carefully revised and the detailed changes are highlighted in \textcolor{blue}{\bf blue color}, along the original text. Our responses to all comments, which were raised by the reviewers, are given in the sequel, using \textcolor{blue}{\bf blue color} as well.

\vi

We hope that those revisions in the manuscript and our responses will be sufficient to make our manuscript suitable for publication in the Journal of Science of Computer Programming (Software Track).

% =============================================
\newpage

\vi
\subsection*{Reviewer \#1 comments followed by the authors' answers:}

\begin{quote}

{\bf Comments to the Authors}

The authors present an extension of the ESBMC tool, to verify CUDA GPU kernels by means of bounded model checking. This is an interesting and relevant topic, since GPUs are used increasingly often to speed up applications significantly. However, programming GPU software can be very tedious and error-prone. A formal verifier of GPU code can make the development process much more efficient.

\begin{quote}
\textcolor{blue}{\textbf{Response:} We really appreciate the reviewer's summary about our work, which fairly describes its main contributions. We also appreciate the reviewer's comments and also his/her effort to provide a relevant discussion, with the goal of improving our paper. As a consequence, we have carefully revised our manuscript, considering all suggestions.}
\end{quote}

{\bf Reviews:}

{\bf (1)} The present submission provides an ok overview of the capabilities of the tool, in my opinion, but the authors never really address the details to fully understand what the tool can and cannot do. For instance, when discussing the experimental evaluation, it is not explained what the source of the bugs is in the various benchmarks.

\begin{quote}
\textcolor{blue}{\textbf{Response:} \textcolor{red}{Lucas: haven't we explained that in {\bf Limitations} subsection on page 6?} \textcolor{green}{Eddie: It seems he meant a detailed explanation about each bug. Do you remember what we did in our former articles? We described each benchmark, which included clarifications about existing bugs... We did not do that here... We only said ``covers basic functions commonly used by real CUDA applications'', but we did not clearly showed what... Maybe a brief explanation of each pat of our test set...} The reviewer wants to know more details about the type of bugs that we have found in the benchmarks. May I ask you to take a look at our journal paper? I think we describe the type of bugs there. }
\end{quote}

  \begin{itemize}
  \item Are there clear differences in the types of bugs that the different tools being compared can detect?
  
  \begin{quote}
  \textcolor{blue}{\textbf{Response:} \textcolor{red}{Hussama: could you answer this question? I believe we need the list of properties that each tool can verify and then we can see the differences between them. We should display such information on the website, in addition to the manuscript.}\textcolor{green}{Eddie: We only do that for our approach (a brief discussion), at the beginning of page 6, above section performance... We need to do that for the other tools... Do we have room for that?}}
  \end{quote}
  
  \item Which types could ESBMC-GPU handle in particular better than the other tools, and why?
  
  \begin{quote}
  \textcolor{blue}{\textbf{Response:} \textcolor{red}{Isabela: you have more experience with the experimental evaluation. What do you think?}\textcolor{green}{Eddie: Maybe, the comment above answers that}}
  \end{quote}
  
  \item To what extent does ESBMC-GPU support CUDA?
  
  \begin{quote}
  \textcolor{blue}{\textbf{Response:} \textcolor{red}{Isabela: would you be able to tell me that?} \textcolor{green}{Eddie: A simple paragraph, just commenting on that and preferably giving a percentage...}}
  \end{quote}
  
  \item Can ESBMC-GPU handle variables in different types of memory? (Shared, global, thread registers?), and inter-warp communication, for instance? 
  
  \begin{quote}
  \textcolor{blue}{\textbf{Response:} \textcolor{red}{Erickson: Could you implement this comment, please?} Lucas: Shall we reference your work about verifying multithread ANSI-C programs? Yes, please. Also reference our journal paper about esbmc-gpu.}
  \end{quote}
  
  \item In that respect, what does COM cover, and what remains to be added?
  
  \begin{quote}
  \textcolor{blue}{\textbf{Response:} \textcolor{red}{Isabela: any insights about coverage?} \textcolor{green}{Eddie: A simple text, just commenting on that and preferably giving a percentage...}}
  \end{quote}
  
  \item Is the memory model for the different types of memory complete?
  
  \begin{quote}
  \textcolor{blue}{\textbf{Response:} \textcolor{red}{Erickson: Could you implement this comment, please?} Lucas: Can you cite our TACAS 2014 paper? I suggest you to read that paper which briefly explains our memory model. Then you can properly reply to the reviewer and cite that paper.\textcolor{green}{Eddie: Maybe he is talking about shared, register, global and local... We need to show that our memory model correctly represent such types...}}
  \end{quote}
  
  \end{itemize}

{\bf (2)} Another paragraph that raises questions is the second on page 3. It seems to say that COM uses pthread functions to simulate GPU kernels. Is that so, and why is it necessary to simulate them? Is this safe to do, because some features in CUDA are not available in C/C++. How do the authors handle those, or do they not?

\begin{quote}
\textcolor{blue}{\textbf{Response:}\textcolor{red}{Erickson: Could you implement this comment, please?} \textcolor{green}{Eddie: We have to comment on kernels, blocks and threads, in order to clarify that, showing that it it feasible ...}}
\end{quote}

{\bf (3)} At some point, the authors mention in passing that Microsoft C++ AMP projects have been used for the evaluation. Does this mean that this language is also supported, or are the benchmarks translated?

\begin{quote}
\textcolor{blue}{\textbf{Response:} Since 2013, ESBMC supports C++. We have added a reference to the paper ``SMT-Based Bounded Model Checking of C++ Programs'', which describes such work.}
\end{quote}

{\bf (4)} Concluding, I would like to see an overall more detailed discussion of the capabilities of the tool.

\begin{quote}
\textcolor{blue}{\textbf{Response:} We have improved Functionalities section, in order to include more details about ESBMC-GPU's capabilities.}
\end{quote}

{\bf (5)} Detailed comments:
\begin{itemize}
\item Page 2
  \begin{itemize}
  \item space-state exploration $\rightarrow$ state-space exploration;
  
    \begin{quote}
    \textcolor{blue}{\textbf{Response:} We have fixed this typo.}
    \end{quote}

  \item ESBMC-GPU builds on top of ESBMC $\rightarrow$ ESBMC-GPU is built on top of ESBMC;
  
    \begin{quote}
    \textcolor{blue}{\textbf{Response:} We have fixed this typo.}
    \end{quote}

  \end{itemize}
\item Page 3
  \begin{itemize}
  \item consists in $\rightarrow$ consists of;
  
    \begin{quote}
    \textcolor{blue}{\textbf{Response:} We have fixed this typo.}
    \end{quote}

  \item Besides, in: remove ``Besides'';
  
    \begin{quote}
    \textcolor{blue}{\textbf{Response:} We have fixed this typo.}
    \end{quote}

  \item reduce state space $\rightarrow$ reduce the state space;
  
    \begin{quote}
    \textcolor{blue}{\textbf{Response:} We have fixed this typo.}
    \end{quote}

  \end{itemize}
\item Page 4
  \begin{itemize}
  \item if multiple threads perform unsynchronized access to shared-memory locations: to the same memory locations, it is crucial to mention that;
  
    \begin{quote}
    \textcolor{blue}{\textbf{Response:} We have implemented as suggested.}
    \end{quote}

  \item has been systematically explored $\rightarrow$ have been systematically explored;
  
    \begin{quote}
    \textcolor{blue}{\textbf{Response:} We have fixed this typo.}
    \end{quote}

  \item Therefore: the next conclusion does not follow logically from the previous statements, so ``therefore'' should be removed here;
  
    \begin{quote}
    \textcolor{blue}{\textbf{Response:} We have fixed this typo.}
    \end{quote}

  \end{itemize}
\item Page 5
  \begin{itemize}
  \item Figure 2: between lines 17 and 18, the contents of {\tt dev\_a} should be copied to {\tt a} with a {\tt cudaMemcpy}, otherwise the subsequent check is quite useless;
  
    \begin{quote}
    \textcolor{blue}{\textbf{Response:} Indeed, there was a code line missing from the figure, we have fixed it now.}
    \end{quote}

  \end{itemize}
\item Page 6
  \begin{itemize}
  \item In addition $\rightarrow$ On the other hand;
  
    \begin{quote}
    \textcolor{blue}{\textbf{Response:} We have implemented as suggested.}
    \end{quote}

  \item Besides, it further simplifies verification models: further compared to what? How does it make it easier?
  
    \begin{quote}
    \textcolor{blue}{\textbf{Response:} We explicitly indicate now that ESBMC-GPU applies MPOR, two-thread analysis, and state hashing in order to further simplify verification models.}
    \end{quote}

  \item if compared: remove ``if'';
  
    \begin{quote}
    \textcolor{blue}{\textbf{Response:} We have fixed this typo.}
    \end{quote}

  \item stream interleaving: what is this?
  
    \begin{quote}
    \textcolor{blue}{\textbf{Response:} We have replaced this term by ``CUDA (parallel) streams and events'' and indicated ``Professional CUDA C Programming'' book as reference.}
    \end{quote}

  \item  to reduce the number of thread interleavings: in what respect? Not reducing the number of threads, right? You already use the 2-thread simplification;
  
    \begin{quote}
    \textcolor{blue}{\textbf{Response:} We intend to use Abstract Interpretation to further reduce the number of thread interleavings (i.e., the model to be verified).}
    \end{quote}

  \end{itemize}
\end{itemize}

\end{quote}

\label{LastPage}

% \printbibliography

\end{document} 
